\section{Second-Order Optimizers} 

	Let's take a look at one of the simplest numerical root-finders.\footnote{Optimization and root finding are very similar, since to optimize $f$ you should first find a root of $f^\prime$.}

  \begin{theorem}[Newton-Raphson Method]
    Given a differentiable function $f: \mathbb{R} \to \mathbb{R}$, the sequence $(x_n)$ defined 
    \begin{equation}
      x_n = x_{n-1} - \frac{f(x_{n-1})}{f^\prime (x_{n-1})}
    \end{equation}
    for any $x_0 \in \mathbb{R}$, converges to a root of $f$. 
  \end{theorem}
  \begin{proof}
    Need to verify this. 
  \end{proof}

  \begin{algo}[Newton-Raphson Method]
    \begin{algorithm}[H]
    \label{alg:newton-raphson}
    \begin{algorithmic}[1]  % The [1] adds line numbers
      % Input/Output documentation
      \Require{Function $f(x)$, initial guess $x_0$}
      \Procedure{NewtonRaphson}{$f, f', x_0, \epsilon, N_{max}$}

      \State $x \gets x_0$ 
      \While{ not converged} 
        \If{$f^\prime (x) = 0$} 
          \State Retry with another $x_0$.  
        \EndIf
        \State $x \gets \frac{f(x)}{f^\prime (x)}$
      \EndWhile
          
      \State \Return $x$ 
      \EndProcedure
    \end{algorithmic}
    \end{algorithm}
  \end{algo}

  \begin{algo}[Secant Method]
    The Secant Method approximates Newton's method by estimating the derivative using two most recent iterates instead of requiring $f'(x)$. 
    \begin{algorithm}[H]
    \label{alg:secant}
    \begin{algorithmic}[1]
      \Require{Function $f(x)$, initial guesses $x_0, x_1$, tolerance $\epsilon$, maximum iterations $N_{max}$}

      \Procedure{Secant}{$f, x_0, x_1, \epsilon, N_{max}$}
        \State $k \gets 0$
        \While{$k < N_{max}$}
          \If{$|f(x_1)| < \epsilon$}
            \State \Return $x_1$
          \EndIf
          \State $x_{new} \gets x_1 - f(x_1)\dfrac{x_1 - x_0}{f(x_1) - f(x_0)}$
          \State $x_0 \gets x_1$
          \State $x_1 \gets x_{new}$
          \State $k \gets k + 1$
        \EndWhile
        \State \Return Failure
      \EndProcedure
    \end{algorithmic}
    \end{algorithm}
  \end{algo}

\subsection{Netwon's Method}

  Newton's method is an iterative algorithm for finding the roots of a differentiable function $F$. An immediate consequence is that given a convex $C^2$ function $f$, we can apply Newton's method to its derivative $f^\prime$ to get the critical points of $f$ (minima, maxima, or saddle points), which is relevant in optimizing $f$. Given a $C^1$ function $f: D \subset \mathbb{R}^n \longrightarrow \mathbb{R}$ and a point $\mathbf{x}_k \in D$, we can compute its linear approximation as 
  \begin{equation}
    f(\mathbf{x}_k + \mathbf{h}) \approx f(\mathbf{x}_k) + D f_{\mathbf{x}_k} \, \mathbf{h} = f(\mathbf{x}_k) + \nabla f(\mathbf{x}_k) \cdot \mathbf{h}
  \end{equation}
  where $D f_{\mathbf{x}_k}$ is the total derivative of $f$ at $\mathbf{x}_k$ and $\mathbf{h}$ is a small $n$-vector. Discretizing this gives us our gradient descent algorithm as 
  \begin{equation}
    \mathbf{x}_{k+1} \gets \mathbf{x}_k - \alpha \, f^\prime(\mathbf{x}_k)
  \end{equation}
  This linear function is unbounded, so we must tune the step size $\alpha$ accordingly. If $\alpha$ is too small, then convergence is slow, and if $\alpha$ is too big, we may overshoot the minimum. Netwon's method automatically tunes this $\alpha$ using the curvature information, i.e. the second derivative. If we take a second degree Taylor approximation 
  \begin{equation}
    f(\mathbf{x}_k + \mathbf{h}) \approx f(\mathbf{x}_k) + D f_{\mathbf{x}_k} \, \mathbf{h} + \mathbf{h}^T \, H f_{\mathbf{x}_k} \, \mathbf{h}
  \end{equation}
  then we are guaranteed that this quadratic approximation of $f$ has a minimum (existence and uniqueness can be proved), and we can calculate it to find our "approximate" minimum of $f$. We simply take the total derivative of this polynomial w.r.t. $\mathbf{h}$ and set it equal to the $n$-dimensional covector $\mathbf{0}$. This is equivalent to setting the gradient as $\mathbf{0}$, so 
  \begin{align*}
    \mathbf{0} & = \nabla_\mathbf{h} \big[ f(\mathbf{x}_k) + D f_{\mathbf{x}_k} \, \mathbf{h} + \mathbf{h}^T \, H f_{\mathbf{x}_k} \, \mathbf{h} \big] (\mathbf{h}) \\
    & = \nabla_\mathbf{h} [ D f_{x_k} \mathbf{h} ] (\mathbf{h}) + \nabla_\mathbf{h} [\mathbf{h}^T \, H f_{\mathbf{x}_k} \, \mathbf{h}] (\mathbf{h}) \\
    & = \nabla_\mathbf{x} f(\mathbf{x}_k) + H f_{\mathbf{x}_k} \, \mathbf{h} \\
    & \implies \mathbf{h} = - [H f_{\mathbf{x}_k}]^{-1} \nabla_\mathbf{x} f(\mathbf{x}_k) 
  \end{align*}
  which results in the iterative update 
  \begin{equation}
    \mathbf{x}_{k+1} \gets \mathbf{x}_k - [H f_{\mathbf{x}_k}]^{-1} \nabla_\mathbf{x} f (\mathbf{x}_k)
  \end{equation}
  Note that we require $\mathbf{f}$ to be convex, so that $H f$ is positive definite. Since $f$ is $C^2$, this implies $H f$ is also symmetric, implying invertibility by the spectral theorem. Note that Newton's method is very expensive, since we require the computation of the gradient, the Hessian, \textit{and} the inverse of the Hessian, making the computational complexity of this algorithm to be $O(n^3)$. We can also add a smaller stepsize $h$ to control stability. 

  \begin{algorithm}
    \caption{Newton's Method}\label{alg:netwons}
    \begin{algorithmic}

    \Require Initial $\mathbf{x}_0$, Stepsize $h$ (optional)

    \For{$t = 0$ to $T$ until convergence}
        \State $g(\mathbf{x}_t) \gets \nabla f(\mathbf{x}_t)$  
        \State $H(\mathbf{x}_t) \gets H f_{\mathbf{x}_t}$ 
        \State $H^{-1} (\mathbf{x}_t) \gets [H(\mathbf{x}_t)]^{-1}$ 
        \State $\mathbf{x}_{t+1} \gets \mathbf{x}_t - h \, H^{-1} (\mathbf{x}_t) \, g(\mathbf{x}_t)$
    \EndFor

    \end{algorithmic}
  \end{algorithm}

	But there is a problem that Netwon's method is prone to saddle points \cite{2014pascanu_saddle-point-nonconvex}. 

\subsection{BFGS}

  Netwon's method is extremely effective for finding the minimum of a convex function, but there are two disadvantages. First, it is sensitive to initial conditions, and second, it is extremely expensive, with a computational complexity of $O(n^3)$ from having to invert the Hessian. An alternative family of optimizers, called \textit{quasi-Newton} methods, try to \textit{approximate} the Hessian (or Jacobian) with $\hat{H} f$, reducing the computational cost without too much loss in convergence rates, and simply use this approximation in the Newton's update: 
  \[\mathbf{x}_{k+1} \gets \mathbf{x}_k - [\hat{H} f_{\mathbf{x}_k}]^{-1} \nabla_\mathbf{x} f (\mathbf{x}_k)\]
  The method of the Hessian approximation varies by algorithm, but the most popular is BFGS. 

  So how do we approximate the Hessian with only the gradient information? With secants. Starting off with $f: \mathbb{R} \longrightarrow \mathbb{R}$, let us assume that we have two points $(x_k, f(x_k))$ and $(x_{k+1}, f(x_{k+1}))$. We can approximate our derivative (gradient in dimension 1) at $x_{k+1}$ using finite differences: 
  \[f^\prime (x_{k+1}) (x_{k+1} - x_k) \approx f(x_{k+1}) - f(x_k)\]
  and doing the same for $f^\prime$ gives us the second derivative approximation: 
  \[f^{\prime\prime} (x_{k+1}) (x_{k+1} - x_k) \approx f^\prime (x_{k+1}) - f^\prime (x_k)\]
  which gives us the update: 
  \[x_{k+1} \gets x_k - \frac{x_{k} - x_{k-1}}{f^\prime (x_k) - f^\prime (x_{k-1})} \, f^\prime (x_k)\]
  This method of approximating Netwon's method in one dimension by replacing the second derivative with its finite difference approximation is called the \textit{secant method}. In multiple dimensions, given two points $\mathbf{x}_k, \mathbf{x}_{k+1}$ with their respective gradients $\nabla f (\mathbf{x}_{k}), \nabla f (\mathbf{x}_{k+1})$, we can approximate the Hessian $\hat{H} f_{\mathbf{x}_{k+1}} \approx D (\nabla f)_{\mathbf{x}_{k+1}}$ (which is the total derivative of the gradient) at $\mathbf{x}_{k+1}$ with the equation
  \[\hat{H} f_{\mathbf{x}_{k+1}} (\mathbf{x}_{k+1} - \mathbf{x}_k) = \nabla_\mathbf{x} f (\mathbf{x}_{k+1}) - \nabla_\mathbf{x} f (\mathbf{x}_k)\]
  This is solving the equation of form $A \mathbf{x} = \mathbf{y}$ for some linear map $A$. Since $\hat{H} f_{\mathbf{x}_{k+1}}$ is a symmetric $n \times n$ matrix with $n (n+1) / 2$ components, there are $n (n+1) / 2$ unknowns with only $n$ equations, making this an underdetermined system. Quasi-Newton methods have to impose additional constraints, with the BFGS choosing the one where we want $\hat{H} f_{\mathbf{x}_{k+1}}$ to be as close as to $\hat{H} f_{\mathbf{x}_{k}}$ at each update $k+1$. Luckily, we can formalize this notion as minimizing the distance between $f_{\mathbf{x}_{k+1}}$ and $\hat{H} f_{\mathbf{x}_{k}}$. Therefore, we wish to find 
  \[\arg \min_{\hat{H} f_{\mathbf{x}_{k+1}}} ||\hat{H} f_{\mathbf{x}_{k+1}} - \hat{H} f_{\mathbf{x}_{k}}||_F,\]
  where $|| \cdot ||_F$ is the Frobenius matrix norm, subject to the restrictions that $\hat{H} f_{\mathbf{x}_{k+1}}$ be positive definite and symmetric and that $\hat{H} f_{\mathbf{x}_{k+1}} (\mathbf{x}_{k+1} - \mathbf{x}_k) = \nabla_\mathbf{x} f (\mathbf{x}_{k+1}) - \nabla_\mathbf{x} f (\mathbf{x}_k)$ is satisfied. Since we have to invert it eventually, we can actually just create an optimization problem that directly computes the inverse. So, we wish to find 
  \[\arg \min_{(\hat{H} f_{\mathbf{x}_{k+1}})^{-1}} ||(\hat{H} f_{\mathbf{x}_{k+1}})^{-1} - (\hat{H} f_{\mathbf{x}_{k}})^{-1} ||_F\]
  subject to the restrictions that 
  \begin{enumerate}
      \item $(\hat{H} f_{\mathbf{x}_{k+1}})^{-1}$ be positive definite and symmetric. It turns out that the positive definiteness restriction also restricts it to be symmetric. 
      \item $\mathbf{x}_{k+1} - \mathbf{x}_k = (\hat{H} f_{\mathbf{x}_{k+1}})^{-1} [\nabla_\mathbf{x} f (\mathbf{x}_{k+1}) - \nabla_\mathbf{x} f (\mathbf{x}_k)]$
  \end{enumerate}
  After some complicated mathematical derivation, which we will not go over here, the problem ends up being equivalent to updating our approximate Hessian at each iteration by adding two symmetric, rank-one matrices $U$ and $V$ scaled by some constant, which can each be computed as an outer product of vectors with itself. 
  \[\hat{H} f_{\mathbf{x}_{k+1}} = \hat{H} f_{\mathbf{x}_{k}} + a U + b V = \hat{H} f_{\mathbf{x}_{k}} + a \mathbf{u} \mathbf{u}^T + b \mathbf{v} \mathbf{v}^T\]
  where $\mathbf{u}$ and $\mathbf{v}$ are linearly independent. This addition of a rank-2 sum of matrices $a U + b V$, known as a rank-2 update, guarantees the "closeness" of $\hat{H} f_{\mathbf{x}_{k+1}}$ to $\hat{H} f_{\mathbf{x}_{k}}$ at each iteration. With this form, we now impose the quasi-Newton condition. Substituting $\Delta \mathbf{x}_k = \mathbf{x}_{k+1} - \mathbf{x}_k$ and $\mathbf{y}_k = \nabla_\mathbf{x} f (\mathbf{x}_{k+1}) - \nabla_\mathbf{x} f (\mathbf{x}_k)$, we have
  \[\hat{H} f_{\mathbf{x}_{k+1}} \Delta \mathbf{x}_k = \hat{H} f_{\mathbf{x}_{k+1}} \Delta \mathbf{x}_k + a \mathbf{u} \mathbf{u}^T \Delta \mathbf{x}_k + b \mathbf{v} \mathbf{v}^T \Delta \mathbf{x}_k = \mathbf{y}_k\]
  A natural choice of vectors turn out to be $\mathbf{u} = \mathbf{y}_k$ and $\mathbf{v} = \hat{H} f_{\mathbf{x}_{k}} \Delta \mathbf{x}_k$, and substituting this and solving gives us the optimal values 
  \[a = \frac{1}{\mathbf{y}_k^T \Delta \mathbf{x}_k}, \;\;\;\;\; b = -\frac{1}{\Delta \mathbf{x}_k^T \hat{H} f_{\mathbf{x}_{k}} \Delta \mathbf{x}_k}\]
  and substituting these values back to the Hessian approximation update gives us the BFGS update: 
  \[\hat{H} f_{\mathbf{x}_{k+1}} = \hat{H} f_{\mathbf{x}_{k}} + \frac{\mathbf{y}_k \mathbf{y}_k^T}{\mathbf{y}_k^T \Delta \mathbf{x}_k} - \frac{\hat{H} f_{\mathbf{x}_{k}} \Delta \mathbf{x}_k \Delta \mathbf{x}_k^T \hat{H} f_{\mathbf{x}_{k}}}{\Delta \mathbf{x}_k^T \hat{H} f_{\mathbf{x}_{k}} \Delta \mathbf{x}_k}\]
  We still need to invert this, and using the \textit{Woodbury formula}
  \[(A + U C V)^{-1} = A^{-1} - A^{-1} U (C^{-1} + V A^{-1} U)^{-1} V A^{-1}\]
  which tells us how to invert the sum of an intertible matrix $A$ and a rank-$k$ correction, we can derive the iterative update of the inverse Hessian as 
  \[(\hat{H} f_{\mathbf{x}_{k+1}})^{-1} = \bigg( I - \frac{\Delta \mathbf{x}_k \mathbf{y}^T}{\mathbf{y}_k^T \Delta \mathbf{x}_k}\bigg) (\hat{H} f_{\mathbf{x}_{k}})^{-1} \bigg( I - \frac{\mathbf{y}_k \Delta \mathbf{x}_k^T}{\mathbf{y}_k^T \Delta \mathbf{x}_k}\bigg) + \frac{\Delta \mathbf{x}_k \Delta \mathbf{x}_k^T}{\mathbf{y}_k^T \Delta \mathbf{x}_k}\]
  Remember that this is the iterative step that we want to actually compute, rather than the ones computing the regular Hessian. The whole point of using the Woodbury formula to derive an inverse update step was to do away with the tedious $O(n^3)$ computations of inverting an $n \times n$ matrix. This rank-2 update also preserves positive-definiteness. 

  Finally, we can choose the initial inverse Hessian approximation $(\hat{H} f_{\mathbf{x}_{k+1}})^{-1}$ to be the identity $I$ or the true inverse Hessian $(H f_{\mathbf{x}_{k+1}})^{-1}$ (computed just once), which would lead to more efficient convergence. The pseudocode for BFGS is a bit too long and confusing to include here, but most of the time, we won't be implementing BFGS by hand; efficient and established BFGS optimizers are already in numerous packages. Like most optimizers, BFGS is not guaranteed to converge to the true global minimum. 

  \begin{algo}[BFGS]
    The BFGS algorithm maintains an approximation $B_k$ to the inverse Hessian, updating it using gradient differences and parameter steps. This avoids explicitly computing or inverting the Hessian at each iteration. 
    \begin{algorithm}[H]
    \label{alg:bfgs}
    \begin{algorithmic}[1]
      \Require{Objective $f(\theta)$, gradient $\nabla f$, initial guess $\theta_0$, tolerance $\epsilon$, maximum iterations $N_{max}$}

      \Procedure{BFGS}{$f, \nabla f, \theta_0$}
        \State $B \gets I$ \Comment{Initialize inverse Hessian approximation}
        \State $\theta \gets \theta_0$
        \State $k \gets 0$
        \While{$k < N_{max}$ and $||\nabla f(\theta)|| > \epsilon$}
          \State $g \gets \nabla f(\theta)$
          \State $d \gets -B g$
          \State $\alpha \gets \text{LineSearch}(f, \theta, d)$ \Comment{Enforce Wolfe conditions}
          \State $\theta_{new} \gets \theta + \alpha d$
          \State $s \gets \theta_{new} - \theta$
          \State $y \gets \nabla f(\theta_{new}) - g$
          \If{$y^T s \leq 0$}
            \State \textbf{break} \Comment{Curvature condition violated}
          \EndIf
          \State $B \gets \Big(I - \frac{s y^T}{y^T s}\Big) B \Big(I - \frac{y s^T}{y^T s}\Big) + \frac{s s^T}{y^T s}$
          \State $\theta \gets \theta_{new}$
          \State $k \gets k + 1$
        \EndWhile
        \State \Return $\theta$
      \EndProcedure
    \end{algorithmic}
    \end{algorithm}
  \end{algo}

  \begin{algo}[Limited-memory BFGS]
    L-BFGS avoids storing the full inverse Hessian by keeping only the last $m$ update pairs $(s_i, y_i)$. 
    At each step, the search direction is computed using a two-loop recursion. 
    This reduces storage from $O(n^2)$ to $O(mn)$ and is widely used in large-scale optimization.
    \begin{algorithm}[H]
    \label{alg:lbfgs}
    \begin{algorithmic}[1]
      \Require{Objective $f(\theta)$, gradient $\nabla f$, initial $\theta_0$, history size $m$, tolerance $\epsilon$, maximum iterations $N_{max}$}

      \Procedure{L-BFGS}{$f, \nabla f, \theta_0, m$}
        \State Initialize $\theta \gets \theta_0$
        \State Initialize empty history lists $S, Y$
        \State $k \gets 0$
        \While{$k < N_{max}$ and $||\nabla f(\theta)|| > \epsilon$}
          \State $g \gets \nabla f(\theta)$
          \State $q \gets g$
          \State Initialize empty list $\alpha$
          \Comment{First loop: backward pass}
          \For{$i = |S|$ down to $1$}
            \State $\rho_i \gets 1 / (y_i^T s_i)$
            \State $\alpha_i \gets \rho_i \, s_i^T q$
            \State $q \gets q - \alpha_i y_i$
          \EndFor
          \State Choose scalar $H_0$ (e.g., $H_0 = \frac{s_{k-1}^T y_{k-1}}{y_{k-1}^T y_{k-1}} I$ if available, else $I$)
          \State $r \gets H_0 q$
          \Comment{Second loop: forward pass}
          \For{$i = 1$ to $|S|$}
            \State $\beta \gets \rho_i \, y_i^T r$
            \State $r \gets r + s_i (\alpha_i - \beta)$
          \EndFor
          \State $d \gets -r$ \Comment{Search direction}
          \State $\alpha \gets \text{LineSearch}(f, \theta, d)$
          \State $\theta_{new} \gets \theta + \alpha d$
          \State $s \gets \theta_{new} - \theta$
          \State $y \gets \nabla f(\theta_{new}) - g$
          \If{$y^T s > 0$} \Comment{Curvature condition}
            \State Append $s, y$ to $S, Y$
            \If{$|S| > m$} 
              \State Remove oldest pair
            \EndIf
          \EndIf
          \State $\theta \gets \theta_{new}$
          \State $k \gets k + 1$
        \EndWhile
        \State \Return $\theta$
      \EndProcedure
    \end{algorithmic}
    \end{algorithm}
  \end{algo}

\subsection{Davidon-Fletcher-Powell (DFP)}

  \begin{algo}[DFP Quasi-Newton Method]
    The DFP method is a quasi-Newton optimization algorithm that maintains an approximation $H_k$ of the inverse Hessian to update parameters efficiently. 
    \begin{algorithm}[H]
    \label{alg:dfp}
    \begin{algorithmic}[1]
      \Require{Objective $f(\theta)$, gradient $\nabla f$, initial guess $\theta_0$, tolerance $\epsilon$}

      \Procedure{DFP}{$f, \nabla f, \theta_0$}
        \State $H \gets I$ \Comment{Initialize inverse Hessian approximation}
        \While{not converged}
          \State $g \gets \nabla f(\theta)$
          \State $d \gets -H g$
          \State $\alpha \gets \text{LineSearch}(f, \theta, d)$
          \State $\theta_{new} \gets \theta + \alpha d$
          \State $s \gets \theta_{new} - \theta$
          \State $y \gets \nabla f(\theta_{new}) - g$
          \State $H \gets H + \frac{ss^T}{s^T y} - \frac{Hy y^T H}{y^T H y}$
          \State $\theta \gets \theta_{new}$
        \EndWhile
        \State \Return $\theta$
      \EndProcedure
    \end{algorithmic}
    \end{algorithm}
  \end{algo}

\subsection{Broyden's Method}

  \begin{algo}[Broyden's Method]
    Broyden's method is a quasi-Newton method that updates an approximation $B_k$ to the Jacobian matrix without computing derivatives directly. 
    \begin{algorithm}[H]
    \label{alg:broyden}
    \begin{algorithmic}[1]
      \Require{Function $F(\theta)$, initial guess $\theta_0$, tolerance $\epsilon$}

      \Procedure{Broyden}{$F, \theta_0$}
        \State $B \gets I$ \Comment{Initial Jacobian approximation}
        \While{not converged}
          \State $\Delta \theta \gets -B^{-1} F(\theta)$
          \State $\theta_{new} \gets \theta + \Delta \theta$
          \State $\Delta F \gets F(\theta_{new}) - F(\theta)$
          \State $B \gets B + \frac{(\Delta F - B \Delta \theta)\Delta \theta^T}{\Delta \theta^T \Delta \theta}$
          \State $\theta \gets \theta_{new}$
        \EndWhile
        \State \Return $\theta$
      \EndProcedure
    \end{algorithmic}
    \end{algorithm}
  \end{algo}

\subsection{Symmetric Rank-1 Update (SR1)}

  \begin{algo}[Symmetric Rank-1 Update]
    The SR1 update is another quasi-Newton method that maintains an approximation $H_k$ of the inverse Hessian, using a symmetric rank-1 correction. 
    \begin{algorithm}[H]
    \label{alg:sr1}
    \begin{algorithmic}[1]
      \Require{Objective $f(\theta)$, gradient $\nabla f$, initial guess $\theta_0$, tolerance $\epsilon$}

      \Procedure{SR1}{$f, \nabla f, \theta_0$}
        \State $H \gets I$
        \While{not converged}
          \State $g \gets \nabla f(\theta)$
          \State $d \gets -H g$
          \State $\alpha \gets \text{LineSearch}(f, \theta, d)$
          \State $\theta_{new} \gets \theta + \alpha d$
          \State $s \gets \theta_{new} - \theta$
          \State $y \gets \nabla f(\theta_{new}) - g$
          \If{$(s - Hy)^T y \neq 0$}
            \State $H \gets H + \dfrac{(s - Hy)(s - Hy)^T}{(s - Hy)^T y}$
          \EndIf
          \State $\theta \gets \theta_{new}$
        \EndWhile
        \State \Return $\theta$
      \EndProcedure
    \end{algorithmic}
    \end{algorithm}
  \end{algo}

